\documentclass[10pt,a4paper]{article}

\usepackage{kotex}
\usepackage{geometry}
\usepackage{setspace}
\usepackage{graphicx}
\usepackage{enumitem}
\usepackage{longtable}
\usepackage{array}

\geometry{
    left=22mm,
    right=22mm,
    top=23mm,
    bottom=23mm
}

\setstretch{1.12}
\setlength{\parindent}{1em}
\setlength{\parskip}{0.25em}
\sloppy

\renewcommand{\abstractname}{요약}
\renewcommand{\refname}{참고문헌}

\begin{document}

\begin{center}
{\LARGE\bfseries TRIPICK: 인구감소지역 관광 활성화를 위한\par}
{\LARGE\bfseries 점수 기반 로컬 여행 동선 추천 서비스 설계 및 구현\par}

\vspace{0.45cm}

이수현$^{1}$, 박준석$^{1}$, 김수린$^{1}$, 이재웅$^{1}$

\vspace{0.2cm}

$^{1}$강원대학교

\vspace{0.2cm}

{\small
suhyun2116@gmail.com, joonsuk0215@gmail.com, surin.kim0714@gmail.com, jaewoong.lee@kangwon.ac.kr
}
\end{center}

\vspace{0.8cm}

\begin{abstract}
최근 국내 인구감소지역은 관광 자원과 로컬 상권을 보유하고 있음에도 불구하고, 낮은 정보 접근성과 부족한 홍보로 인해 여행지로 선택되는 빈도가 낮다. 본 논문에서는 인구감소지역의 숨은 관광지와 로컬 소비 장소를 사용자 조건에 맞게 추천하는 여행 서비스 TRIPICK을 제안한다. TRIPICK은 사용자의 여행 테마, 이동수단, 여행 시간, 동행 유형을 기반으로 장소 후보를 점수화하고, 내부 지역 기여도와 로컬 소비 가능성을 반영하여 동선형 여행 코스를 추천한다. 또한 OpenAI 기반 LLM을 활용하여 추천 이유를 생성하며, LLM 호출 실패 시 룰 기반 추천 이유 생성 방식으로 대체되는 fallback 구조를 적용하였다. 본 시스템은 5개 권역, 89개 지역, 총 4,578개 장소 데이터를 기반으로 구현되었다. 20개 내부 추천 시나리오 기반 검증에서 모든 케이스가 사전에 정의한 적합 기준을 충족하였으며, 실제 추천 응답에 포함된 장소 기준 로컬 소비 장소 포함률은 85.0\%로 나타났다. 본 연구는 단순 관광지 추천을 넘어, 지역 방문 유도와 로컬 소비 촉진을 결합한 지역 활성화형 여행 추천 서비스의 가능성을 제시한다.
\end{abstract}

\noindent
\textbf{주요어:} 인구감소지역, 지역관광, 여행 추천, 동선 추천, 점수 기반 추천, 설명 가능한 추천, LLM, 로컬 소비

\section{서론}

국내 지방 도시와 농어촌 지역은 지속적인 인구 감소와 청년층 유출로 인해 지역 공동체 유지와 경제 활성화에 어려움을 겪고 있다. 특히 인구감소지역은 관광 자원, 전통시장, 로컬 상점, 자연 경관 등 다양한 지역 자산을 보유하고 있음에도 불구하고, 대중적인 여행 플랫폼에서 충분히 노출되지 못하는 경우가 많다. 여행자는 주로 유명 관광지나 대도시 중심의 추천 정보를 접하게 되며, 상대적으로 정보가 부족한 소멸위기지역은 여행 후보지에서 제외되기 쉽다.

기존 여행 추천 서비스는 주로 인기 관광지, 숙소, 맛집, 리뷰 기반 장소 탐색에 초점을 맞추고 있다. 이러한 서비스는 사용자 편의성 측면에서 강점을 가지지만, 인구감소지역의 지역경제 활성화나 로컬 소비 유도라는 사회적 목적을 직접적으로 반영하는 데에는 한계가 있다. 단순히 소멸지역이라는 검색 조건을 추가하는 것만으로는 기존 여행 앱과의 차별성이 충분하지 않으며, 여행자가 해당 지역을 방문해야 하는 이유와 지역에 기여할 수 있는 방식까지 함께 설계할 필요가 있다.

본 연구에서는 이러한 문제를 해결하기 위해 인구감소지역 중심의 로컬 여행 동선 추천 서비스 TRIPICK을 설계하고 구현하였다. TRIPICK은 사용자의 여행 목적과 조건을 반영하여 단일 장소가 아닌 동선 단위의 여행 코스를 추천한다. 또한 추천 결과에는 내부 지역 기여도, 예상 체류시간, 예상 소비액, 로컬 소비 장소 포함 여부를 함께 제공하여 사용자가 여행을 통해 지역 상권과 관광 생태계에 기여할 수 있도록 유도한다.

TRIPICK의 핵심 설계 철학은 LLM이 임의로 장소를 생성하는 방식이 아니라, 실제 TourAPI 및 Supabase DB에 존재하는 장소를 점수 기반으로 선별하고, LLM은 그 결과를 설명하는 역할로 제한하는 것이다. 이를 통해 추천 결과의 안정성과 설명 가능성을 동시에 확보하고자 하였다.

본 논문의 주요 기여는 다음과 같다. 첫째, 인구감소지역을 대상으로 한 로컬 여행 추천 서비스 구조를 제안한다. 둘째, 사용자 조건과 장소 메타데이터를 기반으로 한 점수 기반 추천 로직을 설계하였다. 셋째, 내부 지역 기여도와 로컬 소비 장소 여부를 추천 점수에 반영하여 기존 여행 추천 서비스와의 차별성을 확보하였다. 넷째, LLM 기반 추천 이유 생성과 룰 기반 fallback 구조를 적용하여 추천 결과의 설명 가능성을 높였다. 다섯째, 20개의 내부 추천 시나리오를 통해 추천 결과의 타당성을 검증하였다.

\section{관련 서비스 및 차별점}

현재 국내 여행 관련 서비스로는 여행 일정 관리, 장소 탐색, 리뷰 기반 추천, 숙소 예약, 지도 기반 검색을 제공하는 다양한 플랫폼이 존재한다. 대표적으로 트리플, 대한민국 구석구석, 네이버 지도, 카카오맵, 마이리얼트립 등은 사용자가 여행지를 검색하고, 인기 관광지나 맛집을 탐색하며, 여행 일정을 구성하는 데 도움을 준다. 그러나 이들 서비스는 대부분 사용자의 편의성과 인기 장소 탐색에 초점을 맞추고 있으며, 인구감소지역의 지역 활성화, 로컬 소비 유도, 미션 기반 재방문 유도와 같은 목적을 중심 기능으로 다루지는 않는다.

TRIPICK은 기존 서비스와 비교하여 다음과 같은 차별점을 가진다. 첫째, 서비스 대상 지역을 인구감소지역과 수도권 근교의 상대적 저노출 지역으로 설정하였다. 이는 단순히 유명 여행지를 추천하는 것이 아니라, 상대적으로 방문 기회가 적은 지역을 여행 후보지로 제안한다는 점에서 차별성을 가진다. 둘째, 추천 결과를 단일 장소가 아닌 동선형 코스로 제공한다. 사용자는 하나의 관광지 추천을 받는 것이 아니라, 여행 시간과 이동수단에 맞는 여러 장소의 흐름을 제공받는다. 셋째, 내부 지역 기여도와 로컬 소비 장소를 추천 로직에 반영하였다. 이를 통해 관광지뿐 아니라 전통시장, 로컬 맛집, 카페, 지역 상점 등 지역 소비로 연결될 수 있는 장소가 추천 과정에서 고려된다. 넷째, 미션 및 스탬프 기반 참여 구조를 통해 단발성 방문이 아닌 체류와 재방문을 유도할 수 있도록 설계하였다. 다섯째, 추천 이유 생성에 LLM을 활용하여 사용자가 추천 결과를 이해할 수 있도록 하였다.

\section{서비스 설계}

TRIPICK의 전체 서비스 흐름은 사용자의 여행 조건 입력, 추천 동선 생성, 추천 이유 확인, 지도 기반 동선 확인, 미션 및 스탬프 참여 구조, 추천 히스토리 저장으로 구성된다. 사용자는 먼저 여행 지역을 직접 선택할 수 있으며, 향후에는 사용자 조건을 바탕으로 지역을 추천하는 기능으로 확장할 수 있다. 이후 힐링, 맛집, 뚜벅이, 자연투어, 로컬시장, 지역활성화 추천 등의 테마를 선택하고, 여행 시간, 이동수단, 동행 유형을 입력한다.

본 서비스는 MVP 단계에서 사용자에게 과도한 질문을 하지 않기 위해 핵심 조건을 중심으로 입력 구조를 단순화하였다. 여행 조건은 여행 시간, 이동수단, 동행 유형을 중심으로 구성되며, 사용자가 선택한 테마와 함께 추천 점수 산정에 활용된다. 이러한 방식은 사용자의 명시적 취향 정보를 활용하는 방식으로, 추천 요청 시점의 여행 맥락을 반영할 수 있다.

추천 결과 화면에서는 추천 코스명, 지역 정보, 내부 지역 기여도, 예상 체류시간, 예상 소비액, 로컬 소비 장소 포함 여부, 추천 장소 목록, 추천 이유를 제공한다. 지도 탭에서는 각 추천 장소의 위도와 경도를 기반으로 마커를 표시하고, 사용자가 동선을 직관적으로 확인할 수 있도록 한다. 추천 결과는 사용자의 히스토리에 저장되어 이후 다시 조회할 수 있으며, 향후 개인화 추천의 기초 데이터로 활용될 수 있다.

\section{데이터 구성}

TRIPICK은 한국관광공사 TourAPI를 기반으로 장소 데이터를 수집하고, 서비스 목적에 맞게 지역, 카테고리, 이미지, 로컬 소비 여부, 내부 지역 기여도 등의 속성을 정리하였다. 최종적으로 본 시스템은 5개 권역, 89개 지역, 총 4,578개 장소를 대상으로 구성되었다.

\noindent
\textbf{권역별 데이터 구성은 다음과 같다.}

\begin{itemize}[leftmargin=2em]
\item 충청권: 15개 지역, 865개 장소
\item 강원권: 12개 지역, 710개 장소
\item 수도권 근교: 4개 지역, 253개 장소
\item 전라권: 26개 지역, 1,240개 장소
\item 경상권: 32개 지역, 1,510개 장소
\item 전체: 89개 지역, 4,578개 장소
\end{itemize}

카테고리별 장소 수는 관광지 1,714개, 맛집 1,331개, 액티비티 997개, 문화 364개, 로컬시장 170개, 자연 1개, 카페 1개로 구성되었다. 데이터 출처는 TourAPI 기반 장소 4,574개와 초기 MVP 테스트를 위해 직접 구성한 sample 장소 4개로 이루어져 있다.

카테고리는 TourAPI의 원천 분류를 서비스 목적에 맞게 재분류한 값이다. 자연 경관 장소 중 상당수는 관광지 또는 액티비티 카테고리로 통합되었기 때문에, 자연투어 추천은 단일 category 값이 아니라 theme tag와 장소 설명 정보를 함께 활용하여 판단하였다. 따라서 카테고리 통계에서 자연 항목의 수가 적더라도 자연투어 추천이 불가능하다는 의미는 아니다.

이미지 URL 또는 썸네일 URL이 존재하는 장소는 4,311개이며, 이미지가 없는 장소는 267개로 나타났다. 로컬 소비 장소 여부는 일반 장소 3,053개, 로컬 소비 장소 1,525개로 분류하였다. 로컬 소비 장소는 전통시장, 지역 맛집, 카페, 로컬 상점 등 사용자의 방문이 지역 소비로 이어질 가능성이 높은 장소를 의미한다.

\section{추천 알고리즘 설계}

TRIPICK의 추천 알고리즘은 사용자 조건과 장소 메타데이터를 비교하여 각 장소의 추천 점수를 산정하는 방식으로 설계되었다. 본 시스템은 LLM이 임의로 장소를 생성하거나 선택하지 않도록, 실제 DB에 저장된 장소 후보를 기반으로 점수 기반 추천을 수행한다. 이를 통해 추천 결과의 안정성과 재현 가능성을 확보하였다.

추천 점수는 사용자가 선택한 테마, 이동수단, 동행 유형, 내부 지역 기여도, 로컬 소비 장소 여부를 기준으로 산정된다. 본 논문에서 사용하는 내부 지역 기여도 점수는 공공기관에서 제공하는 공식 지역 활성화 지표가 아니라, 장소의 로컬 소비 가능성, 체류 유도 가능성, 전통시장 및 지역 상권과의 관련성을 바탕으로 서비스 내부에서 정의한 MVP 단계의 추천 보조 점수이다.

\noindent
\textbf{주요 점수 기준은 다음과 같다.}

\begin{itemize}[leftmargin=2em]
\item 테마 일치: 사용자가 선택한 여행 테마와 장소 태그가 일치하는 경우 35점을 부여한다.
\item 이동수단 일치: 선택한 이동수단과 장소 접근성 태그가 일치하는 경우 20점을 부여한다.
\item 이동수단 불일치: 선택한 이동수단과 장소 접근성이 맞지 않는 경우 25점을 감점한다.
\item 동행 유형 일치: 혼자, 친구, 가족, 연인 등 동행 유형과 장소 특성이 일치하는 경우 10점을 부여한다.
\item 내부 지역 기여도: 장소의 지역 활성화 기여 가능성을 반영하기 위해 내부 지역 기여도 점수의 20\%를 추천 점수에 반영한다.
\item 로컬 소비 장소: 맛집, 로컬시장, 지역활성화 테마에서 로컬 소비 장소인 경우 20점을 추가한다.
\item 고기여도 장소: 지역활성화 테마에서 내부 지역 기여도 80점 이상인 장소에는 15점을 추가한다.
\end{itemize}

추천 알고리즘의 흐름은 다음과 같다. 먼저 사용자의 입력 조건을 수집하고, 선택된 권역과 지역에 해당하는 장소 후보를 조회한다. 이후 각 장소의 테마 태그, 이동수단 태그, 동행 유형 태그, 내부 지역 기여도, 로컬 소비 여부를 기준으로 점수를 계산한다. 점수가 높은 장소를 우선 선별하고, 여행 시간에 따라 추천 장소 수를 조절한다. 마지막으로 추천 장소 목록, 지역 정보, 내부 지역 기여도, 예상 체류시간, 추천 이유를 포함한 응답을 생성한다.

본 구조는 LLM만으로 장소를 추천할 때 발생할 수 있는 환각 문제를 줄이고, 실제 DB에 존재하는 장소만을 추천한다는 장점을 가진다. 또한 점수 산정 기준이 명확하기 때문에 추천 결과에 대한 설명과 검증이 가능하다.

\section{LLM 기반 추천 이유 생성}

TRIPICK은 추천 장소 선별 단계에서는 점수 기반 알고리즘을 사용하고, 추천 결과에 대한 설명 단계에서는 LLM을 활용한다. 이는 점수 기반 추천의 안정성과 LLM의 자연어 설명 능력을 결합하기 위한 구조이다.

점수 기반 추천만 사용할 경우 추천 결과의 기준은 명확하지만, 사용자 입장에서는 왜 해당 코스가 추천되었는지 직관적으로 이해하기 어려울 수 있다. 반대로 LLM만으로 추천을 수행할 경우 자연스러운 설명은 가능하지만, 실제 존재하지 않는 장소를 생성하거나 추천 기준이 흔들릴 위험이 있다. 따라서 TRIPICK은 DB 기반 점수 추천을 통해 장소 후보와 동선을 안정적으로 구성하고, LLM은 이미 선별된 추천 결과를 바탕으로 추천 이유를 생성하도록 설계하였다.

현재 구현에서는 OpenAI API를 활용하여 추천 이유를 생성한다. 추천 이유 생성에는 추천 장소 목록, 사용자 선택 테마, 이동수단, 동행 유형, 내부 지역 기여도, 로컬 소비 장소 여부 등의 정보가 활용된다. LLM 호출 시 프롬프트에는 선별된 장소 목록, 사용자 조건, 내부 지역 기여도, 로컬 소비 여부만 제공하며, 새로운 장소를 추가하거나 기존 장소 순서를 변경하지 않도록 제한하였다. LLM 호출이 성공하면 추천 응답에는 LLM 기반 생성 결과임을 나타내는 값이 기록된다. 반면 API 키가 없거나, 네트워크 오류가 발생하거나, 응답 파싱에 실패하는 경우에는 기존 룰 기반 추천 이유 생성 방식으로 대체된다.

이러한 fallback 구조는 외부 LLM API 의존으로 인한 서비스 장애 가능성을 줄이고, 추천 API의 안정성을 확보한다. 즉, LLM은 추천 결과를 대체하는 핵심 결정자가 아니라, 점수 기반 추천 결과를 사용자 친화적으로 설명하는 계층으로 활용된다.

\section{미션 및 스탬프 기반 지역 참여 구조}

본 연구에서는 추천 결과 이후의 지역 참여를 확장하기 위한 구조로 미션 및 스탬프 기능을 설계하였다. MVP 구현에서는 추천 동선 생성, 추천 결과 저장 및 조회, 장소 기반 정보 제공 구조를 중심으로 구현하였으며, 미션 인증과 보상 기능은 향후 서비스 확장 단계에서 연동될 수 있도록 고려하였다.

미션 및 스탬프 구조는 사용자가 추천된 코스 내 장소를 방문하거나, 전통시장 또는 로컬 상점에서 소비하거나, 특정 지역을 재방문하는 방식으로 지역 참여를 수행하도록 유도하는 것을 목표로 한다. 미션 완료 시 스탬프, 쿠폰, 엠블럼 등의 보상을 제공하면 사용자의 참여와 재방문을 유도할 수 있다.

미션 인증 방식은 GPS 기반 위치 인증과 사진 인증을 결합하는 구조를 고려하였다. GPS 기반 인증은 사용자가 실제 장소 근처에 도달했는지를 확인하는 데 활용되며, 사진 인증은 장소 방문 여부를 보조적으로 검증하는 데 활용된다. 또한 향후 사진의 메타데이터를 확인하거나 AI 생성 이미지 여부를 판별하는 방식을 적용하면, 단순 이미지 업로드나 조작된 인증을 줄일 수 있다.

이러한 미션 및 스탬프 구조는 TRIPICK의 핵심 차별점 중 하나이다. 기존 여행 앱이 장소 탐색과 일정 관리 중심이라면, TRIPICK은 지역 방문 이후의 행동까지 설계하여 실제 지역 소비, 체류시간 증가, 재방문 유도와 연결한다. 특히 인구감소지역의 경우 단순 방문 수 증가뿐 아니라 지역 상권과의 접점 형성이 중요하기 때문에, 미션 기반 참여 구조는 지역 활성화 측면에서 의미가 있다.

\section{추천 검증 및 평가}

본 연구에서는 TRIPICK 추천 로직의 타당성을 확인하기 위해 5개 권역을 대상으로 총 20개의 내부 추천 시나리오를 구성하였다. 각 시나리오는 입력 권역, 지역, 테마, 여행 시간, 이동수단, 동행 유형을 다르게 설정하여 실제 추천 응답을 확인하였다. 평가 기준은 테마 적합성, 이동수단 적합성, 로컬 소비 장소 포함 여부, 이미지 제공 여부로 구성하였다.

평가 기준은 다음과 같이 설정하였다. 테마 적합성은 추천 장소 중 70\% 이상이 사용자가 선택한 테마와 관련된 점수 근거를 포함하는 경우 적합으로 판단하였다. 이동수단 적합성은 추천 장소 중 70\% 이상이 선택한 이동수단 조건과 일치하는 경우 적합으로 판단하였다. 로컬 소비 포함 여부는 추천 코스 내에 로컬 소비 장소가 1곳 이상 포함되는지를 기준으로 하였다. 이미지 제공 여부는 전체 장소 데이터가 아니라 20개 내부 추천 시나리오에서 실제 추천 응답에 포함된 장소를 기준으로 산정하였다.

\noindent
\textbf{추천 검증 결과는 다음과 같다.}

\begin{itemize}[leftmargin=2em]
\item 총 테스트 케이스 수: 20건
\item 적합 추천 케이스 수: 20건
\item 유효 추천률: 사전에 정의한 적합 기준 기준 100.0\%
\item 테마 적합률: 사전에 정의한 적합 기준 기준 100.0\%
\item 이동수단 적합률: 사전에 정의한 적합 기준 기준 100.0\%
\item 로컬 소비 장소 포함률: 85.0\%
\item 이미지 제공률: 실제 추천 응답 포함 장소 기준 100.0\%
\end{itemize}

평가 결과, 20개의 내부 추천 시나리오 모두에서 사전에 정의한 테마, 이동수단, 이미지 제공 조건을 충족하였다. 또한 17개 케이스에서 로컬 소비 장소가 포함되어 로컬 소비 장소 포함률은 85.0\%로 나타났다. 이는 TRIPICK의 추천 로직이 정의된 평가 기준 내에서 사용자 조건과 지역 소비 유도 목적을 일정 수준 이상 반영하고 있음을 보여준다. 다만 본 검증은 실제 대규모 사용자 평가가 아니라 내부 시나리오 기반 검증이므로, 향후 실제 사용자 로그와 설문 응답을 활용한 추가 검증이 필요하다.

다만 일부 케이스에서는 태그 기준으로는 적합하나, 실제 여행 감성 측면에서는 사용자가 기대한 테마와 다소 차이가 발생할 가능성이 확인되었다. 예를 들어 자연투어 조건에서 시장이나 식당 중심의 장소가 함께 추천되는 경우가 있었다. 이는 장소 태그의 세분화 부족과 카테고리 중복 문제에서 발생할 수 있으며, 향후 장소 카테고리 정제와 LLM 기반 맥락 검토를 통해 개선할 수 있다.

\section{개인화 추천 고도화 방향}

현재 TRIPICK의 MVP 구현은 추천 요청 시 사용자가 입력하는 테마, 여행 시간, 이동수단, 동행 유형을 명시적 취향 정보로 활용한다. 즉, 현재 단계에서는 사용자의 장기 행동 데이터를 학습하는 개인화 추천 모델이 아니라, 사용자가 매 추천 시점에 제공한 조건을 기반으로 추천 점수를 산정한다.

향후에는 사용자의 여행 기록을 기반으로 개인화 추천을 고도화할 수 있다. 예를 들어 사용자가 저장한 추천 동선, 실제 방문한 장소, 미션 수행 기록, 선호한 테마, 자주 선택한 이동수단, 로컬 소비 장소 방문 여부 등을 누적하면 다음 추천 시 취향 반영 점수에 활용할 수 있다.

향후 개인화 추천 점수는 사용자가 자주 선택한 테마에 대한 가중치, 저장한 동선과 유사한 코스에 대한 선호도, 방문 완료 장소와 유사한 카테고리 선호도, 미션 수행 기록을 통한 로컬 소비 성향, 재방문 지역 또는 미방문 지역에 대한 추천 가중치, 동행 유형별 선호 장소 패턴 등을 포함할 수 있다.

이를 통해 TRIPICK은 단순 조건 기반 추천에서 사용자 행동 기반 개인화 추천으로 확장될 수 있다. 예를 들어 사용자가 전통시장 미션을 자주 완료하고 로컬 맛집을 저장하는 경향이 있다면, 이후 추천에서는 로컬시장과 지역 맛집이 포함된 동선의 점수를 높일 수 있다. 반대로 자연 관광지를 선호하는 사용자는 자연 경관과 산책 코스 중심의 추천을 받을 수 있다.

\section{구현 결과}

TRIPICK의 백엔드는 FastAPI 기반으로 구현되었으며, 추천 API는 사용자 입력 조건을 받아 지역 정보, 추천 동선, 장소 목록, 추천 이유, 내부 지역 기여도, 예상 체류시간, 로컬 소비 장소 포함 여부를 응답한다. 장소 데이터는 Supabase DB에 저장되며, 한국관광공사 TourAPI 기반 데이터와 일부 sample 데이터를 함께 활용하였다.

추천 로직은 서버 코드의 추천 점수 산정 모듈에서 수행되며, 추천 이유 생성 모듈은 OpenAI API 호출을 우선 시도하고 실패 시 룰 기반 템플릿 방식으로 대체된다. 이를 통해 외부 API 장애 상황에서도 추천 응답이 중단되지 않도록 하였다.

추천 결과는 프론트엔드에서 지도 마커로 표시할 수 있도록 각 장소의 위도와 경도 정보를 포함한다. 또한 사용자는 추천 결과를 저장하고, 저장된 동선을 다시 조회할 수 있다. 이 구조는 향후 추천 히스토리 기반 개인화 추천의 기반이 된다.

구현된 추천 API와 저장 및 조회 흐름, TourAPI 기반 데이터 처리, 추천 이유 생성 fallback 동작을 검증하기 위해 총 46개의 단위 및 API 테스트를 수행하였으며, 모든 테스트가 통과하였다.

\section{한계점}

본 연구에는 몇 가지 한계가 존재한다. 첫째, 추천 검증은 20개의 내부 시나리오를 기반으로 수행되었으므로 실제 사용자 전체의 선호를 대표하기에는 한계가 있다. 향후에는 실제 사용자 설문과 서비스 이용 로그를 기반으로 더 많은 테스트 케이스를 확보할 필요가 있다.

둘째, 현재 추천 로직은 장소 태그와 메타데이터를 기반으로 점수를 산정하기 때문에, 태그 품질에 따라 추천 결과의 품질이 달라질 수 있다. 일부 장소는 실제 성격과 태그가 완전히 일치하지 않을 수 있으며, 이로 인해 특정 테마에서 사용자가 기대한 감성과 다소 다른 장소가 추천될 수 있다.

셋째, LLM은 추천 이유 생성에는 활용되지만, 현재 구현에서는 장소 후보를 새로 선택하거나 전체 코스 순서를 재구성하는 핵심 결정자로 사용되지는 않는다. 이는 추천 안정성을 확보하기 위한 설계이지만, 향후에는 점수 기반으로 선별된 후보군 안에서 장소 간 거리, 이동 시간, 분위기 흐름 등을 고려하여 코스 순서를 조정하는 방식으로 LLM 활용 범위를 확장할 수 있다.

넷째, 미션 인증 방식은 GPS와 사진 인증을 중심으로 설계되었으나, 실제 운영 환경에서는 GPS 조작, 이미지 재사용, AI 생성 이미지 업로드 등의 부정 인증 가능성을 추가적으로 고려해야 한다. 이를 해결하기 위해 사진 메타데이터 확인, 촬영 시점 검증, AI 생성 이미지 탐지, 관리자 검수 등의 보완이 필요하다.

\section{결론}

본 논문에서는 인구감소지역의 관광 활성화를 위한 로컬 여행 동선 추천 서비스 TRIPICK을 제안하고 구현하였다. TRIPICK은 사용자의 여행 조건을 기반으로 장소 후보를 점수화하고, 내부 지역 기여도와 로컬 소비 가능성을 반영하여 단일 장소가 아닌 동선형 여행 코스를 추천한다. 또한 LLM 기반 추천 이유 생성과 룰 기반 fallback 구조를 통해 추천 결과의 설명 가능성과 서비스 안정성을 함께 확보하였다.

본 시스템은 5개 권역, 89개 지역, 총 4,578개 장소 데이터를 기반으로 구현되었으며, 20개의 내부 추천 시나리오 검증에서 모든 케이스가 사전에 정의한 적합 기준을 충족하였다. 또한 실제 추천 응답 포함 장소 기준 로컬 소비 장소 포함률은 85.0\%로 확인되었다. 이는 TRIPICK이 사용자 조건에 부합하는 여행 코스를 추천하면서도 지역 소비 유도라는 목적을 정의된 평가 기준 내에서 일정 수준 이상 반영할 수 있음을 보여준다.

향후 연구에서는 실제 사용자 로그와 방문 기록을 기반으로 개인화 추천을 고도화하고, 미션 및 스탬프 수행 데이터를 추천 점수에 반영하는 방식을 확장할 예정이다. 또한 LLM을 활용하여 장소 간 이동 흐름, 여행 분위기, 사용자 선호 맥락을 더 정교하게 반영함으로써 추천 결과의 자연스러움과 만족도를 높일 수 있다. TRIPICK은 인구감소지역의 숨은 관광 자원을 발굴하고, 여행자의 방문을 지역 소비와 재방문으로 연결하는 지역 활성화형 여행 서비스로 발전할 수 있을 것이다.

\section*{부록 A. 추천 시나리오 평가 결과}

\scriptsize
\setlength{\tabcolsep}{3pt}
\renewcommand{\arraystretch}{1.25}

\begin{longtable}{|p{10mm}|p{18mm}|p{31mm}|p{58mm}|p{31mm}|}
\hline
Case & 지역 & 조건 & 대표 추천 장소 & 판정 \\
\hline
TC-01 & 충청권 단양군 & 힐링, 자차, 가족 & 카페산, 만천하스카이워크, 고운골 남한강 갈대숲 & 적합, 로컬 소비 포함, 이미지 제공 \\
\hline
TC-02 & 충청권 공주시 & 맛집, 대중교통, 친구 & 구즉묵보리밥, 그루잠, 고창애풍천장어 & 적합, 로컬 소비 포함, 이미지 제공 \\
\hline
TC-03 & 충청권 태안군 & 자연투어, 자차, 연인 & 갈음이해수욕장, 구례포해수욕장, 기지포해수욕장 & 적합, 로컬 소비 포함, 이미지 제공 \\
\hline
TC-04 & 충청권 부여군 & 로컬시장, 뚜벅이, 혼자 & 부여 중앙시장, 부여장 & 적합, 로컬 소비 포함, 이미지 제공 \\
\hline
TC-05 & 강원권 평창군 & 힐링, 자차, 가족 & 계방산, 고랭길, 고랭지만두마을 & 적합, 이미지 제공, 로컬 소비 미포함 \\
\hline
TC-06 & 강원권 정선군 & 자연투어, 자차, 친구 & 가리왕산, 가리왕산케이블카, 대숲마을 & 적합, 로컬 소비 포함, 이미지 제공 \\
\hline
TC-07 & 강원권 양양군 & 맛집, 대중교통, 연인 & 물치항활어회센터, 거북이서프바 & 적합, 로컬 소비 포함, 이미지 제공 \\
\hline
TC-08 & 강원권 영월군 & 뚜벅이, 뚜벅이, 혼자 & 오그란이 단팥빵, 육정가, 윤스빈 & 적합, 로컬 소비 포함, 이미지 제공 \\
\hline
TC-09 & 수도권 강화군 & 지역활성화, 자차, 가족 & 뜰안에정원, 한옥카페 갤러리 도솔, 꼬막한상 & 적합, 로컬 소비 포함, 이미지 제공 \\
\hline
TC-10 & 수도권 가평군 & 힐링, 자차, 연인 & 가평별묘, 노랑다리미술관, 음악역 1939 & 적합, 이미지 제공, 로컬 소비 미포함 \\
\hline
TC-11 & 수도권 연천군 & 로컬시장, 대중교통, 친구 & 전곡재래시장 & 적합, 로컬 소비 포함, 이미지 제공 \\
\hline
TC-12 & 전라권 담양군 & 맛집, 자차, 친구 & 남도한우, 달빛뜨락, 담양밤부리 & 적합, 로컬 소비 포함, 이미지 제공 \\
\hline
TC-13 & 전라권 보성군 & 힐링, 대중교통, 가족 & 보성 어울마당 야영장, 다향농장, 대계서원 & 적합, 로컬 소비 포함, 이미지 제공 \\
\hline
TC-14 & 전라권 남원시 & 로컬시장, 뚜벅이, 혼자 & 남원 목기공예사, 남원 터미널시장 & 적합, 로컬 소비 포함, 이미지 제공 \\
\hline
TC-15 & 전라권 신안군 & 자연투어, 자차, 연인 & 가거도, 깃대봉, 대광해수욕장 & 적합, 이미지 제공, 로컬 소비 미포함 \\
\hline
TC-16 & 경상권 안동시 & 지역활성화, 대중교통, 가족 & 맘모스베이커리, 맛50년 헛제사밥, 안동 유진찜닭 & 적합, 로컬 소비 포함, 이미지 제공 \\
\hline
TC-17 & 경상권 청도군 & 맛집, 자차, 친구 & 동곡시장, 청도시장, 풍각장 & 적합, 로컬 소비 포함, 이미지 제공 \\
\hline
TC-18 & 경상권 산청군 & 자연투어, 자차, 혼자 & 덕산마을오일장, 산청시장, 산청약초식당 & 적합, 로컬 소비 포함, 이미지 제공, 태그 정제 필요 \\
\hline
TC-19 & 경상권 영주시 & 뚜벅이, 뚜벅이, 연인 & 주막거리캠프, 풍기중앙시장 & 적합, 로컬 소비 포함, 이미지 제공 \\
\hline
TC-20 & 경상권 남해군 & 로컬시장, 자차, 친구 & 남해특산물 지족판매장, 남해시장회센터, 초록스토어 & 적합, 로컬 소비 포함, 이미지 제공 \\
\hline
\end{longtable}

\normalsize
\renewcommand{\arraystretch}{1.0}

\begin{thebibliography}{12}

\bibitem{tourapi}
한국관광공사, 국문 관광정보 서비스 TourAPI, 공공데이터포털, 2026년 6월 접근.

\bibitem{population}
행정안전부, 인구감소지역 지정 현황 및 지원 정책 자료, 2026년 6월 접근.

\bibitem{kto}
한국관광공사, 대한민국 구석구석 서비스, 2026년 6월 접근.

\bibitem{triple}
인터파크트리플, 트리플 여행 서비스, 2026년 6월 접근.

\bibitem{navermap}
네이버, 네이버 지도 서비스, 2026년 6월 접근.

\bibitem{kakaomap}
카카오, 카카오맵 서비스, 2026년 6월 접근.

\bibitem{tripick}
TRIPICK 개발팀, TRIPICK 프로젝트 개발 문서 및 추천 API 명세서, 2026.

\bibitem{survey}
TRIPICK 개발팀, TRIPICK 서비스 기획 설문조사 결과, 2026.

\bibitem{openai}
OpenAI, OpenAI API Documentation, 2026년 6월 접근.

\bibitem{fastapi}
FastAPI, FastAPI Documentation, 2026년 6월 접근.

\bibitem{supabase}
Supabase, Supabase Documentation, 2026년 6월 접근.

\end{thebibliography}

\end{document}

